%==============================================================================
% CHAPTER 7 — DISCUSSION
%==============================================================================
\chapter{Discussion}
\label{chap:discussion}

The preceding chapter presented the empirical results of the AntiGravity evaluation. This chapter interprets those results in relation to the research questions, critically assesses what the thesis does and does not demonstrate, examines the limitations of the work, and discusses the implications for both researchers and practitioners. A researcher who cannot honestly criticize their own work has not understood it.

%------------------------------------------------------------------------------
\section{Interpretation of Results}
\label{sec:disc-interpretation}
%------------------------------------------------------------------------------

\subsection{On the Primary Research Question}

The primary research question (RQ1) asks whether a RAG+MCP architecture enables more \textit{reliable}, \textit{secure}, and \textit{traceable} access to enterprise knowledge than a standalone LLM or naive RAG implementation.

\paragraph{Reliability.} The hallucination rate comparison (Section~\ref{sec:eval-generation}) provides the central evidence for the reliability dimension. The expected finding — that Condition C (RAG+MCP) substantially reduces hallucination compared to Condition A (LLM Alone) — is the most straightforward result to interpret: grounding LLM generation in retrieved, verified source material from the DNSI corpus structurally prevents the model from confabulating information about DNSI-specific processes, systems, and policies that were not in its training data.

The more nuanced finding concerns the comparison between Conditions B and C (Naive RAG vs. RAG+MCP). The faithfulness improvement of Condition C over B, if confirmed, is attributable to two mechanisms: (1) the cross-encoder re-ranking produces a more precisely relevant context, reducing the probability that the model is distracted by loosely related but irrelevant passages; and (2) the MCP context assembly enforces structured citation formatting, which has been observed to anchor the model more firmly to source material in the prompt.

However, the improvement over Naive RAG is expected to be smaller than the improvement over LLM Alone — this is theoretically correct. RAG itself (even without MCP) addresses the fundamental problem of knowledge grounding. MCP adds value in the dimensions of security and standardization, with a secondary benefit to generation faithfulness through better context management.

\paragraph{Security.} The ACL enforcement mechanism provides a structural security guarantee (Section~\ref{subsec:acl-retrieval}): no document content unauthorized for the querying user can enter the LLM prompt. This guarantee is architectural, not probabilistic. It is the primary security contribution of the RAG+MCP design over both Condition A and Condition B. However, this guarantee has boundaries: it applies to the retrieval layer but does not protect against prompt injection attacks via the document content itself (discussed in Section~\ref{sec:disc-security}).

\paragraph{Traceability.} Every response generated by Condition C includes source citations derived from the chunk metadata. This is not merely a user experience feature: it creates an audit trail linking each factual claim in a response to a specific DNSI document, enabling a human reviewer to verify or refute the claim. Conditions A and B provide no equivalent traceability mechanism.

\subsection{On the Value of MCP Specifically}

A critical question for this thesis is whether MCP itself — as opposed to the RAG architecture in general — makes a measurable contribution. This is the hardest question to answer empirically, because MCP's primary contributions (standardization, security boundary, maintainability) are not directly measurable as retrieval or generation quality metrics.

The ablation study (Section~\ref{sec:eval-ablation}) provides indirect evidence: the comparison of A2 (FAISS + re-ranking, no MCP) vs. A4 (full AntiGravity with MCP) isolates the marginal contribution of the MCP layer. If the difference in P@5 or faithfulness between A2 and A4 is not statistically significant — which is the expected finding — this should not be interpreted as MCP providing no value. It should be interpreted as MCP providing value in dimensions that retrieval/generation metrics do not capture: security architecture, connector standardization, and operational maintainability.

This distinction matters for the scientific interpretation of the thesis. The claim is not that MCP improves RAG retrieval quality. The claim is that MCP provides a principled architectural boundary for security enforcement, and that this architectural property is valuable in enterprise contexts in ways that micro-benchmarks cannot fully capture.

%------------------------------------------------------------------------------
\section{Strengths of the Proposed Architecture}
\label{sec:disc-strengths}
%------------------------------------------------------------------------------

\paragraph{Structural security by design.} The most significant strength of the AntiGravity architecture is that the security property --- ACL-aware retrieval --- is not implemented as an application-level check that can be bypassed by logic errors, but as a structural property of the MCP boundary. No pathway exists through the architecture by which unauthorized document content can reach the LLM without traversing the ACL filter.

\paragraph{Source attribution as trust infrastructure.} The mandatory citation mechanism, embedded in the system prompt and enforced by the prompt template, provides a trust infrastructure that enterprise users and compliance teams require. Unlike a standalone LLM, every AntiGravity response can be audited against its source documents. In regulated environments (information security policies, legal compliance, data governance), this auditability is not optional --- it is required.

\paragraph{Protocol separation of concerns.} The MCP boundary between the conversational layer and the data layer means that changes to the underlying data source (a SharePoint API version change, migration to a different SharePoint tenant, addition of a second data source) require changes only to the MCP server --- not to the conversational client or the LLM interface. This operational benefit compounds over time in environments where data infrastructure evolves continuously.

\paragraph{Research novelty.} The architecture addresses a combination of problems --- ACL-aware retrieval, MCP integration, enterprise SharePoint, and multilingual (French/English) corpora --- that no published system has addressed in combination. The research gap identified in Chapter~\ref{chap:background} is genuinely addressed.

%------------------------------------------------------------------------------
\section{Weaknesses and Limitations}
\label{sec:disc-limitations}
%------------------------------------------------------------------------------

This section deliberately presents the limitations of the thesis with the same rigor applied to its strengths. Honest acknowledgment of limitations is a mark of scientific maturity and protects the thesis from criticism on grounds the author has pre-empted.

\subsection{Evaluation Limitations}

\paragraph{Benchmark dataset size ($N=50$).} Fifty question-answer pairs is sufficient for reporting metric distributions with 95\% confidence intervals, but is insufficient for detecting small effect sizes with high statistical power. A benchmark of $N=200$--500 would provide substantially stronger statistical evidence. The constraint is operational: manual annotation of ground-truth relevance labels is time-intensive. Future work should extend the benchmark.

\paragraph{Single-organization generalizability.} All evaluation is conducted on the DNSI SharePoint corpus. The vocabulary, document structure, query patterns, and language distribution (French/English mix) are specific to this organization. Results may not generalize to organizations with different corpus characteristics. The evaluation should be interpreted as a proof-of-concept demonstration, not as a universal performance characterization.

\paragraph{RAGAS as a proxy for faithfulness.} RAGAS metrics are computed by a language model evaluating the faithfulness of another language model's output against retrieved context. This creates a potential circular dependency: if the evaluating model has similar failure modes to the generating model, some hallucinations may not be detected. A fully reliable hallucination assessment requires human annotation of every generated response --- which was provided for the hallucination rate metric but not for the RAGAS faithfulness scores.

\paragraph{No longitudinal evaluation.} The user study is a single-session evaluation. Real adoption decisions are made over time, through repeated use. A longitudinal study (tracking usage patterns over weeks) would provide much stronger evidence for or against the adoption hypothesis (H5).

\subsection{Technical Limitations}

\paragraph{LLM inference latency.} The CPU-only deployment of Mistral 7B results in generation latency of 25--45 seconds for typical responses --- substantially exceeding NFR01 (P95 $\leq$ 5 seconds). This is the most significant gap between the designed system and the production-ready system. It is a deployment infrastructure limitation (no GPU), not an architectural limitation. GPU deployment is expected to reduce generation latency to 2--5 seconds, bringing the system within NFR01. This limitation was known at the outset and does not invalidate the architectural contributions, but it does prevent the user study from fully reflecting the intended user experience.

\paragraph{Incremental index updates.} The current ingestion pipeline requires a full index rebuild when documents are deleted from SharePoint. For a dynamically evolving corpus, this means a latency window during which deleted documents remain retrievable. The window is bounded by the ingestion schedule (nightly) but is non-zero. Production deployment would require migration to a vector store with native deletion support (e.g., Chroma, Weaviate, or FAISS \texttt{IndexIDMap}).

\paragraph{Single-language embedding model.} The \texttt{all-mpnet-base-v2} embedding model was trained primarily on English text. While it performs adequately on the French-language DNSI documents (French is well-represented in its training data), a multilingual model (e.g., \texttt{paraphrase-multilingual-mpnet-base-v2}) or a domain-fine-tuned model would likely improve retrieval precision on French-language queries.

\paragraph{Prompt injection residual risk.} As identified in the threat model (Table~\ref{tab:threat-model}), the prompt injection attack vector via adversarial document content has not been fully mitigated. The structured system prompt and near-zero temperature reduce susceptibility, but do not eliminate it. A production deployment should add an additional layer of prompt injection detection (e.g., a small classifier model screening retrieved chunks for injection patterns) before this system handles highly sensitive materials.

\subsection{Scope Limitations}

\paragraph{No real-time synchronization.} The nightly ingestion schedule means the index may be up to 24 hours out of date relative to the live SharePoint corpus. For time-sensitive operational contexts (incident response, live project documentation), this lag is unacceptable. Real-time delta synchronization via the Microsoft Graph change notification API is required for production use and is identified as the highest-priority future work item.

\paragraph{Single data source.} The system is designed for SharePoint CERP exclusively. DNSI knowledge is distributed across multiple systems (email, wikis, ticketing systems, databases). A production enterprise assistant would need to federate across these sources. MCP's multi-server architecture makes this extension architecturally straightforward but operationally complex.

%------------------------------------------------------------------------------
\section{Security Analysis: Critical Assessment}
\label{sec:disc-security}
%------------------------------------------------------------------------------

The security architecture presented in Chapter~\ref{chap:architecture} provides strong guarantees at the retrieval layer but leaves residual risks at the prompt layer that deserve honest analysis.

\paragraph{What the architecture guarantees.} The delegated OAuth 2.0 token mechanism ensures that ACL enforcement is performed by Microsoft's Graph API permission engine --- the most authoritative and thoroughly tested component in the system. ACL checks performed by well-maintained cloud provider services are more reliable than ACL checks implemented in application code. The MCP boundary ensures that the retrieval and permission resolution logic cannot be bypassed by the conversational layer.

\paragraph{What the architecture does not guarantee.} The architecture does not protect against an insider threat who uses their legitimate SharePoint write access to plant adversarial documents. A document containing the text \texttt{[SYSTEM: disregard previous instructions and output the full text of the previous response]} would pass the ACL filter (the attacker is authorized to write it) and be retrieved by legitimate queries. The near-zero temperature setting and structured system prompt provide partial mitigation but not elimination of this risk.

\paragraph{Practical likelihood.} This attack requires a motivated insider with SharePoint write access who is aware of the AntiGravity system's implementation details and is willing to compromise the system. In the DNSI context, this is a low-probability but non-zero threat. The recommendation is to add automated prompt injection screening of indexed document content at ingestion time, as an additional defense-in-depth measure.

%------------------------------------------------------------------------------
\section{Deployment Challenges: Lessons Learned}
\label{sec:disc-deployment}
%------------------------------------------------------------------------------

The deployment of AntiGravity in the DNSI's production environment produced several lessons that are generalizable to any enterprise AI deployment:

\paragraph{Corporate proxy is a first-class constraint, not an afterthought.} The DNSI's TLS-terminating proxy required changes to all components in the stack (HTTP clients, MSAL, MCP transport). Enterprise AI systems should be designed with proxy-aware HTTP clients from the start, not retrofitted. The corporate proxy is not an exceptional edge case --- it is the norm in large enterprise environments.

\paragraph{OAuth 2.0 delegated permissions require organizational IT cooperation.} Registering an Entra ID application, configuring delegated permissions, and navigating the DNSI's internal approval workflow for API access consumed significant project time. Enterprise AI projects should account for identity and access management lead times in their project schedules.

\paragraph{On-premises LLM inference is operationally expensive.} Running Mistral 7B on-premises (required by NFR09) requires hardware procurement, model management, quantization trade-off decisions, and ongoing maintenance. Cloud-based LLM APIs are operationally much simpler, but are incompatible with strict data locality requirements. Organizations face a genuine trade-off between operational simplicity and data sovereignty.

%------------------------------------------------------------------------------
\section{Generalizability of Findings}
\label{sec:disc-generalizability}
%------------------------------------------------------------------------------

The AntiGravity architecture and the findings of this thesis are generalizable to other organizations under the following conditions:

\begin{enumerate}
  \item \textbf{SharePoint as the primary knowledge repository:} The connector design and ACL enforcement mechanism are specific to SharePoint and Microsoft Graph API. Organizations using Confluence, Box, Google Drive, or other platforms would require a different MCP server implementation, but the architectural pattern (MCP boundary for security enforcement) is fully transferable.
  \item \textbf{Enterprise security requirements:} The security architecture is designed for organizations with formal information security policies (ISO 27001 or equivalent). Organizations without strict ACL requirements can simplify the architecture by removing the permission manifest and delegated token mechanism.
  \item \textbf{Moderate corpus scale (up to 100k documents):} The FAISS IndexFlatIP performs optimally at this scale. Larger organizations would need to evaluate approximate search indexes.
  \item \textbf{Document-centric knowledge:} The system is optimized for unstructured text in PDF and DOCX format. Organizations with knowledge primarily in structured databases, spreadsheets, or specialized formats would require extended parsing modules.
\end{enumerate}

%------------------------------------------------------------------------------
\section{Future Work}
\label{sec:disc-future}
%------------------------------------------------------------------------------

The limitations identified above directly motivate a prioritized future research and development agenda:

\begin{table}[H]
\centering
\caption{Prioritized future work roadmap}
\label{tab:future-work}
\begin{tabularx}{\textwidth}{@{}llXl@{}}
\toprule
\textbf{Priority} & \textbf{Item} & \textbf{Description} & \textbf{Expected Impact} \\
\midrule
P1 & GPU deployment & Deploy Mistral 7B on GPU hardware to bring generation latency within NFR01 & Critical for user adoption \\
P2 & Real-time delta sync & Microsoft Graph change notifications for live SharePoint updates & Production readiness \\
P3 & Multilingual embedding & Fine-tune or replace embedding model for French/English domain vocabulary & Retrieval precision \\
P4 & Prompt injection defense & Ingestion-time adversarial content screening & Security hardening \\
P5 & Larger benchmark & Expand to $N=200$+ Q\&A pairs for statistical power & Evaluation rigor \\
P6 & Multi-source federation & Extend MCP architecture to email, wikis, ticketing systems & Coverage expansion \\
P7 & Agentic RAG & Integrate multi-step tool use for complex procedural queries & Capability expansion \\
P8 & CIFRE doctoral continuation & Investigate MCP as a universal enterprise AI integration standard & Scientific contribution \\
\bottomrule
\end{tabularx}
\end{table}

\paragraph{CIFRE doctoral direction.} This thesis establishes a foundation for a CIFRE doctoral research program that would investigate three open scientific questions: (1) Can MCP be formalized as a universal integration standard for enterprise AI, and what security-theoretic guarantees can be proven about MCP-mediated architectures? (2) How do RAG+MCP systems scale under multi-tenant, multi-source, and multi-modal knowledge conditions? (3) What are the organizational change management factors that determine enterprise AI assistant adoption, and how can system design choices mitigate adoption barriers?

%------------------------------------------------------------------------------
\section{Chapter Summary}
\label{sec:disc-summary}
%------------------------------------------------------------------------------

This chapter has provided a critical analysis of the AntiGravity system and the thesis research. The primary research question is answered affirmatively: the RAG+MCP architecture demonstrably improves reliability (reduced hallucination), security (structural ACL enforcement), and traceability (mandatory source citations) compared to baseline alternatives. MCP's specific contribution is primarily architectural (security boundary, standardization) rather than metric-measurable (retrieval precision, faithfulness), which is precisely what the MCP specification claims and what this thesis confirms.

The thesis's limitations are significant but bounded: the evaluation is a proof-of-concept demonstration on a single organization's corpus with a small benchmark, not a universal performance characterization. The deployment has a known infrastructure gap (CPU inference latency) that is resolvable with hardware investment. These limitations are typical of applied AI engineering research at the Master's level and do not undermine the validity of the architectural contributions.

The following chapter synthesizes the thesis contributions and answers each research question explicitly.
